\section{Preliminaries on hidden-action principal-agent problems}\label{sec:preliminaries}
% In the following Section~\ref{sec:classical_problem}, we formally introduce the principal-agent problems that we study in this work.
%
% Then, in Section~\ref{sec:learning_problem} we introduce the learning task which we tackle in the rest of the paper.


% \subsection{Hidden-action principal-agent problem}\label{sec:classical_problem}


An instance of the \emph{hidden-action principal-agent problem} is characterized by a tuple $\left(\mathcal{A},\Omega\right)$, where $\mathcal{A}$ is a finite set of $n \coloneqq |\mathcal{A}|$ actions available to the agent, while $\Omega$ is a finite set of $m \coloneqq |\Omega|$ possible outcomes.
%
Each agent's action $a \in \mathcal{A}$ determines a probability distribution $F_a \in \Delta_{\Omega}$ over outcomes, and it results in a cost $ c_a \in [0, 1]$ for the agent.\footnote{Given a finite set $X$, we denote by $\Delta_X$ the set of all the probability distributions over the elements of $X$.}
%
We denote by $F_{a,\omega}$ the probability with which action $a$ results in outcome $\omega \in \Omega$, as prescribed by $F_a$.
%
Thus, it must be the case that $\sum_{\omega \in \Omega}F_{a,\omega}=1$ for all $a \in \mathcal{A}$.
%
Each outcome $\omega \in \Omega$ is characterized by a reward $r_\omega \in [0,1]$ for the principal.
%
Thus, when the agent selects action $a \in \mathcal{A}$, the principal's expected
reward is $\sum_{\omega \in \Omega}F_{a,\omega}r_{\omega}$. 


The principal commits to an outcome-dependent payment scheme with the goal of steering the agent towards desirable actions.
%
Such a payment scheme is called \emph{contract} and it is encoded by a vector $p \in \mathbb{R}^{m}_{+}$ defining a payment $p_\omega \geq 0$ from the principal to the agent for each outcome $\omega \in \Omega$.\footnote{As it is customary in contract theory~\citep{carroll2015robustness}, in this work we assume that the agent has \emph{limited liability}, meaning that the payments can only be from the principal to the agent, and \emph{not viceversa}.}
%
% In this work, we assume that the contracts representing valid commitments for the principal lie in a polytope $\mathcal{P} \subset \mathbb{R}^{m}_{+}$, and we denote by $h_{\preg} \in \mathbb{N}$ the number of linear inequalities defining such a polytope. 
%
%The principal’s goal is to commit to a contract that maximizes their expected utility, as formally described in the following. A contract $p \in \mathcal{P} \subset \mathbb{R}^{m}_{+}$ specifies the payments from the principal to the agent for each possible realized outcome $\omega \in \Omega$, \emph{i.e.}, $p_\omega$. In the following we assume $ \mathcal{P} \subset \mathbb{R}^{m}_{+}$
%to be a closed convex polytope defined by the intersection of $h(\preg)$ half spaces.
%
% \alb{Ho messo pedice perchè non è una funzione. Valutare se dire qualcosa dei vincoli di non-negatività.} \mat{dire perchè, sembra un assunzione totalmente assurda. magari vediamo cosa dice MJ. capisco che è presto per dire cosi a caso non si può fare learning}
%
%
% As it is customary in principal-agent problems, we assume the payments to be non-negative, \emph{i.e.}, $p_\omega \ge 0$ for each $\omega \in \Omega$.
%
% Consequently, when the agent selects action $a \in \mathcal{A}$ the principal's expected utility is given by $u(p) \coloneqq  \sum_{\omega \in \Omega}F_{a,\omega} \left( r_\omega - p_{\omega} \right)$, while the agent's expected utility is given by $\sum_{\omega \in \Omega}F_{a,\omega}p_{\omega} - c_a$.
%
Given a contract $p \in \mathbb{R}^{m}_{+}$, the agent plays a \emph{best-response} action that is: {(i)} \emph{incentive compatible} (IC), \emph{i.e.}, it maximizes their expected utility; and {(ii)} \emph{individually rational} (IR), \emph{i.e.}, it has non-negative expected utility.
%
%
%\begin{itemize}
%    \item \emph{incentive compatible} (IC), \emph{i.e.}, it maximizes their expected utility; and
%    \item \emph{individually rational} (IR), \emph{i.e.}, it has non-negative expected utility.
%    % (if there is no IR action, then the agent abstains from playing so as to preserve the status quo).
%\end{itemize}
%
We assume w.l.o.g. that there always exists an action $a \in \mathcal{A}$ with $c_a=0$.
%
Such an assumption ensures that there is an action providing the agent with positive utility.
%
This guarantees that any IC action is also IR and allows us to focus w.l.o.g. on IC only.
%
%To focus only on IC only actions we rely on the folllowing assumptions:
%\begin{assumption}
%There always exists an action $a \in \mathcal{A}$ such that $c_a=0$.
%\end{assumption}
%Intuitively, the following assumption ensures that there exists an action providing the agent with a non-negative utility, thus ensuring IR of any IC action. 



% As a result, 
%
Whenever the principal commits to $p \in \mathbb{R}^{m}_{+}$, the agent's expected utility by playing an action $a \in \mathcal{A}$ is equal to $\sum_{\omega \in \Omega}F_{a,\omega}p_{\omega} - c_a$, where the first term is the expected payment from the principal to the agent when selecting action $a$.
%
Then, the set $A(p) \subseteq \mathcal{A}$ of agent's best-response actions in a given contract $p \in \mathbb{R}^{m}_{+}$ is formally defined as follows:
%\begin{equation*}%\label{eq:best_response}
    $A(p) \coloneqq  \arg\max_{a \in \mathcal{A}}  \sum_{\omega \in \Omega}F_{a,\omega}p_{\omega} - c_a .$
%\end{equation*}
%
Given an action $a \in \mathcal{A}$, we denote with $\mathcal{P}_{a} \subseteq \mathbb{R}^{m}_{+}$ the \emph{best-response set} of action $a$, which is the set of all the contracts that induce action $a$ as agent's best response.
%
Formally, $\mathcal{P}_{a} \coloneqq \left \{ p \in \mathbb{R}^{m}_{+} \mid  a \in A(p) \right \}$.
%\begin{equation*}
%$$\mathcal{P}_{a} \coloneqq \left \{ p \in \mathbb{R}^{m}_{+} \mid \sum_{\omega \in \Omega}F_{a,\omega}p_{\omega} - c_a \ge  \sum_{\omega \in \Omega}F_{a',\omega}p_{\omega} - c_{a'} \,\,\, \forall a' \in \mathcal{A} \right \}.$$
%\end{equation*}
%
%We say that an action $a \in \mathcal{A}$ is \emph{implementable} if $\mathcal{P}_{a}$ is non-empty.\alb{Meglio usare \emph{feasible}?}\bac{bisognerebbe anche agggiunger che il contratto p implenta l'azione a.}
%
\begin{comment}
We also introduce the definition of the spaces of contracts in which the expected utility of the agent by playing action $a_i$ is greater or equal to the one achieved by playing action $a_j$, as formally stated in the following: 
\begin{equation*}
\mathcal{P}_{ij}\coloneqq \left \{ p \in \mathbb{R}^m \, | \, \sum_{\omega \in \Omega}F_{a_i ,\omega}p_{\omega} - c_{a_i}  \ge  \sum_{\omega \in \Omega}F_{a_j ,\omega}p_{\omega} - c_{a_j} \right \}.
\end{equation*} 
As a consequence, for each action $a_i \in \mathcal{A}$, we have that $\mathcal{P}_{a_i} \coloneqq \bigcap_{j \in [n]} \mathcal{P}_{ij}$. 
As a final step we introduce the formal defintion of the hyperplanes separating the region $\mathcal{H}_{ij}$ form $\mathcal{H}_{ij}$, for each couple of actions $a_i,a_j \in \mathcal{A}$. Formally we have:
\begin{equation}
H_{ij} \coloneqq \left \{ x \in \mathbb{R}^m \, | \, \sum_{\omega \in \Omega}F_{a_i ,\omega}x_{\omega} - c_{a_i}  \ge  \sum_{\omega \in \Omega}F_{a_j ,\omega}x_{\omega} - c_{a_j} \right \}.
\end{equation}
\end{comment}


As it is customary in the literature (see, \emph{e.g.}, \citep{dutting2019simple}), we assume that the agent breaks ties in favor of the principal when having multiple best responses available.
%
We let $a^\star(p) \in A(p)$ be the action played by the agent in a given $p \in \mathbb{R}^{m}_{+}$, which is an action $a \in A(p)$ maximizing the principal's expected utility $\sum_{\omega \in \Omega}F_{a,\omega} \left( r_\omega - p_{\omega} \right)$.
%
For ease of notation, we introduce the function $u: \mathbb{R}^{m}_{+} \to \mathbb{R}$ to encode the principal's expected utility under all the possible contracts; formally, the function $u$ is defined so that $u(p) = \sum_{\omega \in \Omega}F_{a^\star(p),\omega} \left( r_\omega - p_{\omega} \right)$ for every $p \in \mathbb{R}^{m}_{+}$.


In the classical (single-round) hidden-action principal-agent problem, the goal of the principal is find a contract $p \in \mathbb{R}^{m}_{+}$ that maximizes the expected utility $u(p)$.
%\footnote{The single-round version of the problem can be solved in polynomial time with an LP~\citep{dutting2021complexity}.}
%
By letting $\mathsf{OPT} := \max_{p \in \mathbb{R}^{m}_{+}} u(p)$, we say that a contract $p \in \mathbb{R}^{m}_{+}$ is \emph{optimal} if $u(p) = \mathsf{OPT}$.
%
Moreover, given an additive approximation error $\rho > 0$, we say that $p$ is \emph{$\rho$-optimal} whenever $u(p) \geq \mathsf{OPT} -\rho$.




\section{Learning optimal contracts}\label{sec:learning_problem}
We study settings in which the principal and the agent interact over multiple rounds, with each round involving a repetition of the same instance of hidden-action principal-agent problem.
%
The principal has no knowledge about the agent, and their goal is to learn an optimal contract.
%from experience.
%
% \emph{multi-round} version of the classical hidden-action principal-agent problem in which the principal has no knowledge about the agent, and, thus, the former has to learn an $\epsilon$-optimal contract by repeatedly interacting with the latter.
At each round, the principal-agent interaction goes as follows: (i) The principal commits to a contract $p \in \mathbb{R}^{m}_{+}$. (ii) The agent plays $a^\star(p)$, which is \emph{not} observed by the principal. (iii) The principal observes the outcome $\omega \sim F_{a^\star(p)}$ realized by the agent's action.
\begin{comment}
\begin{enumerate}
	\item[(i)] The principal commits to a contract $p \in \mathbb{R}^{m}_{+}$.
	%
	\item[(ii)] The agent plays action $a^\star(p)$, which is \emph{not} observed by the principal.
	%
	\item [(iii)] The principal observes the outcome $\omega \sim F_{a^\star(p)}$ realized by the agent's action.
\end{enumerate}
\end{comment}
%
The principal only knows the set $\Omega$ of possible outcomes and their associated rewards $r_\omega$, while they do \emph{not} know anything about agent's actions $\mathcal{A}$, their associated distributions $F_a$, and their costs $c_a$. 
%
%Differently from classical contract design problems, we consider a repeated interaction between the principal and the agent. Specifically, at each round the principal commits to a contract $p \in \mathcal{P}$, receiving as feedback the outcome sampled from the probability distribution $F_{a(p)}$.
%In the following we assume that at each iteration the principal to commit to a contract $p \in \mathcal{P}$, observing only observing the realized outcome. Formally we assume that at each iteration the principal-agent interaction unfold as follows: (i) the principal commits to a contract $p \in \preg$; (ii) The agent plays action $a(p)$ (unkown to the principal) according to Equation \ref{eq:best_response}; (iii) The principal observes the outcome $\omega \sim F_{a(p)}$.
%
% \alb{Per fare LP finale servono anche i reward. C'è un modo per farlo senza? In quel caso saremmo in grado di generalizzare il setting di Jordan nel caso regret minimization.}\bac{si ma poi cosa osserveresti ad ogni iterazione? $r_\omega$ ? Se cosi fosse non è come osservare $\omega$? } \mat{la risposta penso che sia: imprarare il reward è triviale, basta osservare una volta che quell'outcome è uscito. Quindi per semplicità fingiamo di conoscerlo. Ma check di cosa fa MJ}\bac{MJ osserva outcome} 


The goal is to design algorithms for the principal that prescribe how to select a contract at each round in order to learn an optimal contract.
%
Ideally, we would like an algorithm that, given an additive approximation error $\rho > 0$ and a probability $\delta \in (0,1)$, is guaranteed to identify a $\rho$-optimal contract with probability at least $1-\delta$ by using the minimum possible number of rounds.\footnote{Such a learning goal can be seen as instantiating the PAC-learning framework into principal-agent settings.}
%
The number of rounds required by such an algorithm can be seen as the \emph{sample complexity} of learning optimal contracts in hidden-action principal-agent problems (see~\citep{zhu2022sample}). 
%
%directly connected to what is known as the \emph{sample complexity} of learning optimal contracts (see also~\citep{zhu2022sample}). 


The following theorem shows that the goal defined above is too demanding.
%
\begin{restatable}{theorem}{spacehardness}\label{thm_hardness_space} 
	For any number of rounds $N \in \mathbb{N}$, there is no algorithm that is guaranteed to find a $\kappa$-optimal contract with probability greater than or equal to $1 - \delta$ by using less than $N$ rounds, where $\kappa, \delta > 0$ are some suitable absolute constants.
\end{restatable}
%
Theorem~\ref{thm_hardness_space} follows by observing that there exist instances in which an approximately-optimal contract may prescribe a large payment on some outcome.
%
Thus, for any algorithm and number of rounds $N$, the payments used up to round $N$ may \emph{not} be large enough to learn such an approximately-optimal contract. To circumvent Theorem~\ref{thm_hardness_space}, in the rest of the paper we focus on designing algorithms for the problem introduced in the following Definition~\ref{def:learning_prob}.
%
Such a problem relaxes the one introduced above by asking to find a contract whose seller's expected utility is ``approximately'' close to that of a utility-maximizing contract among those with  \emph{bounded} payments. %above by a given value $B \geq 1$.
%
% Formally:
%
\begin{definition}[Learning an optimal bounded contract]\label{def:learning_prob}
	The problem of \emph{learning an optimal bounded contract} reads as follows: Given a bound $B \geq 1$ on payments, an additive approximation error $\rho > 0$, and a probability $\delta \in (0,1)$, find a contract $p \in [0,B]^m$ such that the following holds:
	\[
		\mathbb{P} \left\{  u(p) \geq \max_{p' \in [0,B]^m} u(p') - \rho \right\} \geq 1-\delta.
	\]
\end{definition}
%
Let us remark that the problem introduced in Definition~\ref{def:learning_prob} does \emph{not} substantially hinder the generality of the problem of learning an optimal contract.
%
Indeed, since a contract achieving principal's expected utility $\mathsf{OPT} $ can be found by means of linear programming~\citep{dutting2021complexity}, it is always the case that an optimal contract uses payments which can be bounded above in terms of the number of bits required to represent the probability distributions over outcomes.
%
Moreover, all the previous works that study the sample complexity of learning optimal contracts focus on settings with bounded payments; see, \emph{e.g.},~\citep{zhu2022sample}.
%
The latter only considers contracts in $[0,1]^m$, while we address the more general case in which the payments are bounded above by a given $B \geq 1$.
% and our analysis unveils the dependency of the sample complexity on $B$.


%In such a multi-round setting, the goal of the principal is to learn an $\epsilon$-optimal contract as fast as possible, \emph{i.e.}, by minimizing the total number of rounds $T \in \mathbb{N}$ of the repeated principal-agent interaction.
%%
%Finding such a minimum value of $T$ is also known as the problem of establishing the \emph{sample complexity of online contract design} in the literature \alb{Citare}.
%%
%Moreover, such a problem is intimately connected with doing regret minimization in online hidden-action principal-agent problems.
%%
%Our results also extend to such a regret-minimization framework, as we extensively discuss in Section~\ref{sec:regret}.\bac{il T grande di solito è il time-horizon nei problemi bandit di regret minimization magari se si puo usare una altra lettera}
%
%
%The principal's goal is to learn an (approximately) optimal contract while minimizing the number of (sub-optimal) iterations. In the following we denote with $\mathsf{OPT}$ the highest expected utility that the principal can achieve. Furthermore, given $\epsilon > 0$, we say that a contract $p \in \mathcal{P}$ is $\epsilon$-optimal if the principal's expected utility in $p$ satisfies $u(p) \ge \textnormal{OPT} - \epsilon$. In other words, an $\epsilon$-optimal contract is one where the principal's expected utility is within $\epsilon$ of the optimal value $\textnormal{OPT}$. 

